% Audit de veracite des citations de rapport de jury - 2026-09-20
% Bloc 4 : citation jury 2025 remplacee par le verbatim exact du
% Rapport_Jury_2025 (l.1024). Les mentions << 85 % des candidats >> et
% << gain de l'amplificateur differentiel >> sont introuvables dans le corpus.
\expandafter\ifx\csname CAPETINCLUDEMODE\endcsname\relax
\documentclass[11pt,a4paper]{article}
% preamble.tex — packages et commandes communes a toutes les fiches CAPET SII IE
\input{glyphtounicode}
\pdfgentounicode=1
\usepackage{cmap}
\usepackage[T1]{fontenc}
\usepackage[utf8]{inputenc}
\usepackage[french]{babel}
% babel[french] active a \begin{document} des raccourcis de ponctuation (:;!?)
% qui cassent la syntaxe to[...] de circuitikz (erreur "+ or - expected") ;
% on les desactive, ce qui ne change que l'espacement fin avant ces signes.
\AtBeginDocument{\shorthandoff{:;!?}}
\usepackage[a4paper, top=1.6cm, bottom=1.6cm, left=1.4cm, right=1.4cm]{geometry}
\setlength{\headheight}{14.5pt}
\usepackage{amsmath, amssymb, amsfonts}
\usepackage{siunitx}
\sisetup{locale=FR, output-decimal-marker={,}, inter-unit-product=\ensuremath{{}\cdot{}},
         group-minimum-digits=4}
% Unites non fournies par siunitx, declarees ici une fois pour toutes (fiches + exercices).
% Les gardes \ifdefined evitent tout conflit avec une redeclaration locale restee dans un
% fichier : ne pas retirer les gardes, ne pas redeclarer ces unites dans les fichiers.
\DeclareSIUnit{\jour}{j}
\ifdefined\var\else\DeclareSIUnit{\var}{var}\fi   % puissance reactive
\ifdefined\vad\else\DeclareSIUnit{\vad}{VAD}\fi   % puissance deformante
\ifdefined\tr\else\DeclareSIUnit{\tr}{tr}\fi      % tour (tr/min)
\usepackage{array, booktabs, multirow, tabularx}
\usepackage{xcolor}
\usepackage[most, breakable]{tcolorbox}
\usepackage{multicol}
\usepackage{enumitem}
\usepackage{mathtools}
\usepackage{lmodern}
% Refonte visuelle 2026 (direction "moderne coloré") : Lato pour le texte,
% maths sans-serif assorties via newtxsf (charge apres amsmath/amssymb).
\usepackage[default]{lato}
\usepackage{newtxsf}
\usepackage{microtype}
\usepackage{fancyhdr}
\usepackage{tikz}
\usetikzlibrary{positioning,arrows.meta}
\usepackage{circuitikz}
\usepackage{lastpage}

% ── Couleurs par domaine (palette modernisée 2026) ────────────────────
\definecolor{colElec}{RGB}{29, 78, 216}        % Électrotechnique — bleu vif
\definecolor{colElecLight}{RGB}{219, 234, 254} % fond bleu clair
\definecolor{colAuto}{RGB}{21, 128, 61}        % Automatique — vert
\definecolor{colMeca}{RGB}{194, 65, 12}        % Mécanique — orange
\definecolor{colTherm}{RGB}{185, 28, 28}       % Thermique — rouge
\definecolor{colQT}{RGB}{232, 244, 233}        % (hérité, plus utilisé par questiontype)
\definecolor{colJury}{RGB}{255, 248, 225}      % fond jury (crème)
\definecolor{colJuryFrame}{RGB}{202, 138, 4}   % bordure/bandeau jury (ambre)
\definecolor{colGrayLight}{RGB}{245, 245, 247} % fond gris neutre
\definecolor{colSeparator}{RGB}{180, 180, 180} % gris séparateur
\definecolor{colExo}{RGB}{8, 145, 178}          % Banque d'exercices — cyan (distinct des fiches)
% ── Encre et gris de la refonte ───────────────────────────────────────
\definecolor{colInk}{RGB}{23, 32, 42}          % texte courant (quasi-noir)
\definecolor{colSlate}{RGB}{83, 98, 118}       % gris ardoise (méta, en-têtes)
\definecolor{colRule}{RGB}{203, 213, 225}      % filets clairs
\definecolor{colJuryBg}{RGB}{255, 248, 225}    % fond crème jury
\definecolor{colPresence}{RGB}{83, 98, 118}    % bloc 5 (présence dans les sujets)
\AtBeginDocument{\color{colInk}}

% ── Style encadré exercice + corrigé (banque d'exercices) ─────────────
\tcbset{
  exoBox/.style={
    breakable,
    enhanced,
    colback=colExo!8!white,
    colframe=colExo,
    colbacktitle=colExo,
    coltitle=white,
    fonttitle=\bfseries\small,
    left=4pt, right=4pt, top=3pt, bottom=3pt,
    before skip=5pt, after skip=5pt,
  }
}

% ── Bascule enonce / corrige (banque d'exercices) ─────────────────────
% \CAPETCORRIGE est defini en ligne de commande (pdflatex "\def\CAPETCORRIGE{1}\input{...}")
% pour compiler la variante corriges d'un meme fichier source ; par defaut (non defini),
% la compilation produit la variante enonces. \CAPETSHOW{enonce}{corrige} choisit la branche.
%
% TROIS MODES (refonte 2026-09-20, manuel de revision) :
%   - aucun drapeau   -> imprime l'enonce            (variante enonces, INCHANGE)
%   - \CAPETCORRIGE   -> imprime le corrige          (variante corriges, INCHANGE)
%   - \CAPETMANUEL    -> imprime l'enonce ET met le corrige de cote pour l'annexe
%
% Le mode manuel stocke chaque corrige dans une macro (\csgdef), sans ecriture
% fichier ni re-tokenisation : le contenu (circuitikz, accents, %) traverse intact.
% L'annexe les rejoue ensuite dans l'ordre via \CAPETANNEXE.
\usepackage{etoolbox}

\newcounter{CAPETexo}          % numero continu d'exercice dans le manuel
\setcounter{CAPETexo}{0}
\newcounter{CAPETannexe}       % compteur de relecture en annexe

% Titre courant de la notion en cours (pose par \CAPETNOTION), sert d'intertitre annexe
\newcommand{\CAPETnotioncode}{}
\newcommand{\CAPETnotiontitre}{}

% \CAPETNOTION{code}{titre} : ouvre une notion dans le manuel (memorise le libelle,
% et marque le rang du premier exercice de la notion pour l'intertitre d'annexe).
\newcommand{\CAPETNOTION}[2]{%
  \renewcommand{\CAPETnotioncode}{#1}%
  \renewcommand{\CAPETnotiontitre}{#2}%
  \csgdef{CAPETnotionstart\the\numexpr\value{CAPETexo}+1\relax}{#1 -- #2}%
}

% ── \degree robuste ───────────────────────────────────────────────────
% Certains fichiers d'exercices utilisent \degree. Ni le preambule ni siunitx
% ne le definissent, et \providecommand local ne suffit pas : sous hyperref,
% \degree existe deja comme commande factice (\ERROR), donc \providecommand
% s'abstient et l'appel echoue a la compilation du manuel (erreur invisible
% en compilation isolee). On le (re)definit donc inconditionnellement ici,
% apres le chargement des packages, pour tout le corpus.
\AtBeginDocument{\providecommand{\degree}{}\renewcommand{\degree}{\ensuremath{^\circ}}}

% Libelle decoratif : les fichiers d'exercices utilisent aussi \CAPETSHOW dans
% leur \fancyhead et leur bloc de titre, pour afficher "(enonces)" ou "(corriges)".
% Ces appels-la ne sont PAS des exercices : en mode manuel ils ne doivent ni
% consommer un numero, ni alimenter l'annexe. On les reconnait a leur premier
% argument, strictement egal au libelle ci-dessous.
\newcommand{\CAPETLIBELLEENONCES}{(\'enonc\'es)}

\ifdefined\CAPETMANUEL
  % --- Mode manuel : enonce en place, corrige differe vers l'annexe ---
  \newcommand{\CAPETSHOW}[2]{%
    \def\CAPETargun{#1}%
    \ifx\CAPETargun\CAPETLIBELLEENONCES
      #1%            appel decoratif : rendu tel quel, hors numerotation
    \else
      \CAPETSHOWEXO{#1}{#2}%
    \fi
  }
  \newcommand{\CAPETSHOWEXO}[2]{%
    \stepcounter{CAPETexo}%
    \csgdef{CAPETcorrige\theCAPETexo}{#2}%
    \edef\CAPETtmpnum{\theCAPETexo}%
    #1%
    \par\vspace{2pt}%
    {\footnotesize\itshape\color{colExo}%
     \hyperref[capet:cor:\CAPETtmpnum]{Corrige de l'exercice \CAPETtmpnum\ $\rightarrow$ p.~\pageref{capet:cor:\CAPETtmpnum}}%
     \label{capet:eno:\CAPETtmpnum}}%
  }
  % \CAPETANNEXE : rejoue tous les corriges captures, dans l'ordre, avec intertitres
  \newcommand{\CAPETANNEXE}{%
    \setcounter{CAPETannexe}{0}%
    \loop\ifnum\value{CAPETannexe}<\value{CAPETexo}%
      \stepcounter{CAPETannexe}%
      \edef\CAPETcurnum{\theCAPETannexe}%
      \ifcsname CAPETnotionstart\CAPETcurnum\endcsname
        \par\medskip
        {\normalsize\bfseries\color{colExo}\csuse{CAPETnotionstart\CAPETcurnum}}%
        \par\nobreak\smallskip
      \fi
      \begin{tcolorbox}[exoBox, breakable,
        title={Corrige de l'exercice \CAPETcurnum\
               \hfill\normalfont\footnotesize
               \hyperref[capet:eno:\CAPETcurnum]{enonce p.~\pageref{capet:eno:\CAPETcurnum}}}]%
      \label{capet:cor:\CAPETcurnum}%
      \csuse{CAPETcorrige\CAPETcurnum}%
      \end{tcolorbox}%
    \repeat
  }
\else
  % --- Modes historiques, comportement strictement inchange ---
  \newcommand{\CAPETSHOW}[2]{\ifdefined\CAPETCORRIGE #2\else #1\fi}
  \newcommand{\CAPETANNEXE}{}
\fi

% ── Difficulté 1 à 3 (étoiles pleines, amssymb) ───────────────────────
\newcommand{\difficulte}[1]{%
  \ifcase#1\relax
  \or $\bigstar$\;niveau 1/3%
  \or $\bigstar\bigstar$\;niveau 2/3%
  \or $\bigstar\bigstar\bigstar$\;niveau 3/3%
  \fi
}

% ── Style encadré extrait de document technique (banque d'exercices) ──
\tcbset{
  docTech/.style={
    breakable,
    enhanced,
    colback=colGrayLight,
    colframe=colSeparator,
    fonttitle=\bfseries\small\itshape,
    coltitle=black,
    colbacktitle=colGrayLight,
    left=4pt, right=4pt, top=3pt, bottom=3pt,
    before skip=4pt, after skip=4pt,
  }
}

% ── Style tcolorbox de section principale ────────────────────────────
% Refonte 2026 : bandeau titre coloré arrondi + fond très pâle, sans cadre.
\tcbset{
  ficheSection/.style={
    breakable, enhanced jigsaw,
    frame hidden, boxrule=0pt, arc=1.8mm,
    colback=#1!4!white,
    colbacktitle=#1,
    coltitle=white,
    fonttitle=\sffamily\bfseries,
    left=7pt, right=7pt, top=5pt, bottom=5pt,
    toptitle=3pt, bottomtitle=3pt,
    before skip=10pt, after skip=8pt,
  }
}

% ── Style encadré question type ──────────────────────────────────────
% Carte blanche, bandeau sombre, liseré très clair.
\tcbset{
  questiontype/.style={
    breakable, enhanced jigsaw,
    colback=white,
    frame hidden, boxrule=0pt, arc=1.8mm,
    colbacktitle=colElec!78!black,
    coltitle=white,
    fonttitle=\sffamily\bfseries\small,
    left=7pt, right=7pt, top=4pt, bottom=4pt,
    toptitle=2.5pt, bottomtitle=2.5pt,
    before skip=8pt, after skip=8pt,
    borderline={0.6pt}{0pt}{colElec!12!white},
  }
}

% ── Style encadré jury ───────────────────────────────────────────────
% Bandeau ambre, fond crème, sans cadre.
\tcbset{
  jury/.style={
    breakable, enhanced jigsaw,
    colback=colJuryBg,
    frame hidden, boxrule=0pt, arc=1.8mm,
    colbacktitle=colJuryFrame,
    coltitle=white,
    fonttitle=\sffamily\bfseries,
    left=7pt, right=7pt, top=4pt, bottom=4pt,
    toptitle=3pt, bottomtitle=3pt,
    before skip=10pt, after skip=8pt,
  }
}

% ── En-tête / pied de page communs (refonte 2026) ─────────────────────
% #1 code fiche, #2 titre court, #3 couleur domaine
\newcommand{\ficheHeaders}[3]{%
  \fancyhead[L]{\small\sffamily{\bfseries\textcolor{#3}{#1}}\hspace{0.7em}\textcolor{colSlate}{#2}}%
  \fancyhead[R]{\small\sffamily\textcolor{colSlate}{CAPET SII IE\, --- \,Session 2026}}%
  \fancyfoot[C]{\small\sffamily\textcolor{colSlate}{\thepage\,/\,\pageref{LastPage}}}%
  \fancyhead[C]{}\fancyfoot[L]{}\fancyfoot[R]{}%
}

% ── Bloc en-tête de fiche (BLOC 0, refonte 2026) ──────────────────────
% #1 code  #2 titre  #3 couleur domaine  #4 domaine  #5 ligne présence  #6 sous-thèmes
\newcommand{\ficheEntete}[6]{%
  \begin{tcolorbox}[enhanced, breakable,
    frame hidden, boxrule=0pt, arc=2.5mm,
    interior style={left color=#3, right color=#3!75!black},
    left=9pt, right=9pt, top=7pt, bottom=7pt,
    before skip=0pt, after skip=6pt]
    {\color{white}\sffamily\bfseries\LARGE #1\quad #2}\\[4pt]
    {\color{#3!22!white}\sffamily\small \MakeUppercase{#4}\quad\textbar\quad #5}\\[1pt]
    {\color{#3!35!white}\sffamily\footnotesize Sous-thèmes : #6}
  \end{tcolorbox}%
}

% ── Mise en page ─────────────────────────────────────────────────────
\pagestyle{fancy}
\fancyhf{}
\renewcommand{\headrulewidth}{0.5pt}
\renewcommand{\footrulewidth}{0pt}
\AtBeginDocument{%
  \renewcommand{\headrule}{{\color{colRule}\hrule width\headwidth height 0.5pt}}%
}

\setlist[itemize]{leftmargin=1.2em, itemsep=1pt, topsep=2pt, parsep=0pt}
\setlist[enumerate]{leftmargin=1.2em, itemsep=3pt, topsep=2pt, parsep=0pt}

\setlength{\parskip}{2pt}
\setlength{\columnsep}{10pt}

\def\CAPETBEGINFILE{\begin{document}}
\def\CAPETENDFILE{\end{document}}
\else
\def\CAPETBEGINFILE{}
\def\CAPETENDFILE{}
\fi

\ficheHeaders{T23}{Capteurs et conditionnement}{colElec}

\CAPETBEGINFILE

% -----------------------------------------------------------------------
%  BLOC EN-TETE
% -----------------------------------------------------------------------

\ficheEntete{T23}{Capteurs et Conditionnement du signal}{colElec}%
  {électronique - Instrumentation industrielle}%
  {\textbf{Présence :} \textbf{9/13} sujets (2013-2026) - thème \emph{très fréquent}}%
  {PT100 - thermocouple - capteur inductif - jauge de contrainte - pont de Wheatstone - signal 4-20\,mA - montages AOP de conditionnement - mise à l'échelle CAN}

% -----------------------------------------------------------------------
%  1. RAPPELS DE COURS
% -----------------------------------------------------------------------

\begin{tcolorbox}[
  ficheSection=colElec,
  title={1. Rappels de cours},
]

% -- 1.1 Taxonomie capteurs ---------------------------------------------
\textbf{1.1 Classification des capteurs}

\smallskip
\begin{center}
\renewcommand{\arraystretch}{1.25}
\begin{tabular}{lll}
\toprule
\textbf{Critère} & \textbf{Actif} & \textbf{Passif} \\
\midrule
Principe & Génère une FEM & Varie une impédance \\
Exemples & Thermocouple, piézoélectrique & PT100, jauge, potentiomètre \\
Alimentation & Non (source autonome) & Oui (pont, courant) \\
\bottomrule
\end{tabular}
\end{center}

\medskip\hrule\vspace{4pt}

% -- 1.2 PT100 ---------------------------------------------------------
\textbf{1.2 Capteur PT100 (thermorésistance)}

\[
  R_T = R_0(1 + \alpha\,T)
  \qquad R_0 = 100\,\Omega \text{ à } 0^\circ\mathrm{C}
  \qquad \alpha = 3{,}85 \times 10^{-3}\,^\circ\mathrm{C}^{-1}
\]
\[
  \frac{\Delta R_T}{\Delta T} = R_0\,\alpha = 0{,}385\,\Omega/^\circ\mathrm{C}
\]
Capteur \textbf{passif} ; plage $-200^\circ$C à $+850^\circ$C. Nécessite un pont d'alimentation.

\medskip\hrule\vspace{4pt}

% -- 1.3 Thermocouple --------------------------------------------------
\textbf{1.3 Thermocouple (effet Seebeck)}

Capteur \textbf{actif} : génère une FEM $e_{AB}$ de quelques mV entre deux métaux différents.
\begin{center}
\renewcommand{\arraystretch}{1.2}
\begin{tabular}{lll}
\toprule
\textbf{Type} & \textbf{Plage} & \textbf{Tolérance} \\
\midrule
J (Fe/Constantan) & $-210^\circ$C à $+1200^\circ$C & $\pm 2{,}5^\circ$C \\
K (Chromel/Alumel) & $-10^\circ$C à $+1100^\circ$C  & $\pm 2{,}5^\circ$C \\
\bottomrule
\end{tabular}
\end{center}
Nécessite un transmetteur mV → 4-20\,mA ou un ampli d'instrumentation.

\medskip\hrule\vspace{4pt}

% -- 1.4 Pont de Wheatstone --------------------------------------------
\textbf{1.4 Pont de Wheatstone - jauge de contrainte}

\[
  V_{out} \approx \frac{V_{cc}}{4} \cdot \frac{\Delta R}{R_0} \quad (\Delta R \ll R_0)
\]

\smallskip
\begin{center}
\begin{circuitikz}[european, scale=0.85]
  \draw (0,2) node[above]{$+V_{cc}$} to[R, l=$R_1$] (2,1)
        to[R, l=$R_2$] (0,0) node[below]{0 V};
  \draw (0,2) to[R, l_=$R_4$] (-2,1)
        to[R, l_=$R_3{+}\Delta R$] (0,0);
  \draw (2,1) -- (2.8,1) node[right]{$V_+$};
  \draw (-2,1) -- (-2.8,1) node[left]{$V_-$};
\end{circuitikz}
\end{center}

Avantage : compense les dérives thermiques (4 jauges identiques). L'ampli différentiel amplifie $V_{out} = V_+ - V_-$ jusqu'à la plage du CAN.

\medskip\hrule\vspace{4pt}

% -- 1.5 Signal 4-20 mA ------------------------------------------------
\textbf{1.5 Signal normalisé industriel 4-20\,mA}

\[
  x = x_{min} + \frac{I_s - 4}{16}\,(x_{max} - x_{min})
\]
\begin{itemize}
  \item $I_s = 4\,\mathrm{mA}$ - grandeur minimale (pas zéro !)
  \item $I_s = 20\,\mathrm{mA}$ - grandeur maximale
  \item Immunité aux perturbations CEM sur grande longueur de câble (2 fils)
\end{itemize}

\medskip\hrule\vspace{4pt}

% -- 1.6 Conditionnement AOP -------------------------------------------
\textbf{1.6 Montages AOP courants en conditionnement}

\begin{center}
\renewcommand{\arraystretch}{1.3}
\begin{tabular}{lll}
\toprule
\textbf{Montage} & \textbf{Gain} & \textbf{Usage} \\
\midrule
Inverseur & $-R_f/R_e$ & Amplification + inversion \\
Non-inverseur & $1 + R_f/R_e$ & Amplification sans inversion \\
Soustracteur (diff.) & $(R_2/R_1)(V^+ - V^-)$ & Pont de Wheatstone, PT100 \\
Suiveur & $1$ & Adaptation d'impédance \\
\bottomrule
\end{tabular}
\end{center}

\end{tcolorbox}

% -----------------------------------------------------------------------
%  2. FORMULAIRE SYNTHETIQUE
% -----------------------------------------------------------------------

\begin{tcolorbox}[
  ficheSection=colElec,
  title={2. Formulaire synthétique},
]
\small
\begin{multicols}{2}
\renewcommand{\arraystretch}{1.45}
\begin{tabular}{@{}p{4.8cm}ll@{}}
\toprule
\textbf{Formule} & \textbf{Usage} & \textbf{U.} \\
\midrule
\multicolumn{3}{@{}l}{\footnotesize\textit{PT100}} \\
$R_T = R_0(1 + \alpha T)$ & Résistance & $\Omega$ \\
$\Delta R_T/\Delta T = R_0\alpha = 0{,}385$ & Sensibilité & $\Omega/^\circ$C \\
\midrule
\multicolumn{3}{@{}l}{\footnotesize\textit{Pont Wheatstone}} \\
$V_{out} \approx (V_{cc}/4) \cdot (\Delta R/R_0)$ & Déséquilibre & V \\
$V_{CAN} = K_a \cdot V_{out}$ & Après ampli & V \\
\midrule
\multicolumn{3}{@{}l}{\footnotesize\textit{Signal 4-20\,mA}} \\
$x = x_{min} + (I_s-4)/16 \cdot (x_{max}-x_{min})$ & Mise à l'éch. & - \\
\bottomrule
\end{tabular}

\columnbreak

\renewcommand{\arraystretch}{1.45}
\begin{tabular}{@{}p{4.8cm}ll@{}}
\toprule
\textbf{Formule} & \textbf{Usage} & \textbf{U.} \\
\midrule
\multicolumn{3}{@{}l}{\footnotesize\textit{Conditionnement}} \\
$V_s = -(R_f/R_e) V_e$ & Inverseur & V \\
$V_s = (1+R_f/R_e) V_e$ & Non-inv. & V \\
$V_s = (R_2/R_1)(V^+-V^-)$ & Soustr. & V \\
\midrule
\multicolumn{3}{@{}l}{\footnotesize\textit{Capteur inductif}} \\
$f = N \cdot Z \cdot n / 60$ & Fréquence signal & Hz \\
\midrule
\multicolumn{3}{@{}l}{\footnotesize\textit{Données PT100}} \\
$R_0 = 100\,\Omega$ ; $\alpha = 3{,}85\times10^{-3}\,^\circ$C$^{-1}$ & - & - \\
Plage : $-200^\circ$C à $+850^\circ$C & - & - \\
\bottomrule
\end{tabular}
\end{multicols}
\normalsize
\end{tcolorbox}

% -----------------------------------------------------------------------
%  3. QUESTIONS TYPES
% -----------------------------------------------------------------------

\begin{tcolorbox}[
  ficheSection=colElec,
  title={3. Questions types - extraites des sujets réels},
]

% -- QT1 -------------------------------------------------------------

\begin{tcolorbox}[questiontype,
  title={QT1 - Identification du capteur et justification du conditionnement \hfill\normalfont\footnotesize E2 BAC3 2026 (Q1-Q2)}
]
\textbf{Énoncé.}
Capteur Jumo (référence 370) placé dans une conduite. Identifier le type, le principe de mesure, sa nature active/passive. Justifier le besoin d'un conditionneur pour $T \in [-5\,;\,90^\circ\mathrm{C}]$, $V_{CAN} \in [0\,;\,3] \,\mathrm{V}$. Montrer que $c = 0{,}0316\,\mathrm{V/^\circ C}$.

\tcbline

\textbf{Corrigé.}

\textbf{Étape 1 — Identification.} Type : PT100 (thermorésistance) - transducteur résistif, capteur \textbf{passif}.

\textbf{Étape 2 — Justification du conditionneur.} Plage capteur : variation de résistance correspondant à $95^\circ$C. Tension de sortie directe du capteur : trop faible et décalée - le conditionneur adapte la plage.

\textbf{Étape 3 — Calcul du coefficient.}
\[
  c = \dfrac{V_{max} - V_{min}}{T_{max} - T_{min}} = \dfrac{3 - 0}{90 - (-5)} = \dfrac{3}{95} \approx \mathbf{0{,}0316\,\mathrm{V/^\circ C}}
\]

\textbf{Conclusion.} Ce coefficient $c$ garantit que toute la plage utile de température ($-5$ à $90^\circ$C) est traduite sur l'intégralité de la plage d'entrée du CAN ($0$ à $3$\,V), maximisant ainsi la résolution effective de la mesure.
\end{tcolorbox}

\vspace{4pt}

% -- QT2 -------------------------------------------------------------

\begin{tcolorbox}[questiontype,
  title={QT2 - Pont de Wheatstone - expression de la tension de déséquilibre \hfill\normalfont\footnotesize E2 Commune 2025}
]
\textbf{Énoncé.}
Un pont de Wheatstone alimente 4 jauges de contrainte ($R_0 = 120\,\Omega$, $V_{cc} = 1\,\mathrm{V}$). Une jauge subit $\Delta R = 0{,}0153\,\Omega$. (a) Expliquer l'intérêt du pont par rapport à une mesure directe. (b) Calculer $V_{out}$.

\tcbline

\textbf{Corrigé.}

\textbf{Étape 1 (a) — Intérêt du pont.} Le pont compense les dérives thermiques : si toutes les jauges varient de la même façon avec la température, la variation s'annule. Seule la déformation mécanique appliquée à une jauge crée un déséquilibre.

\textbf{Étape 2 (b) — Calcul.} Approximation petits déséquilibres ($\Delta R \ll R_0$) :
\[
  V_{out} \approx \frac{V_{cc}}{4} \cdot \frac{\Delta R}{R_0} = \frac{1}{4} \cdot \frac{0{,}0153}{120} \approx \mathbf{31{,}9\,\mu\mathrm{V}}
\]
Gain ampli nécessaire pour 5\,V en CAN : $K_a = 5 / (31{,}9 \times 10^{-6}) \approx 1{,}56 \times 10^5$.

\textbf{Conclusion.} Un gain de l'ordre de $1{,}56 \times 10^5$ est très élevé pour un AOP classique : il impose l'usage d'un amplificateur d'instrumentation dédié (offre de dérive et de bruit maîtrisés), pas d'un simple montage soustracteur à AOP standard.
\end{tcolorbox}

\vspace{4pt}

% -- QT3 -------------------------------------------------------------

\begin{tcolorbox}[questiontype,
  title={QT3 - Capteur inductif - calcul de fréquence et choix \hfill\normalfont\footnotesize E1 2020}
]
\textbf{Énoncé.}
Capteur inductif détectant les dents d'un disque. Vitesse $n = 240\,\mathrm{tr/min}$, nombre de dents $Z = 24$. Calculer la fréquence maximale du signal. Choisir le capteur adapté parmi la liste.

\tcbline

\textbf{Corrigé.}

\textbf{Étape 1 — Calcul de fréquence.}
\[
  f = \frac{N \cdot Z \cdot n}{60} = \frac{1 \times 24 \times 240}{60} = \mathbf{96\,\mathrm{Hz}}
\]

\textbf{Étape 2 — Choix du capteur.} Le capteur choisi doit avoir une fréquence de commutation maximale $f_{max} \geq 96\,\mathrm{Hz}$.

\textbf{Conclusion.} Comparer $f_{max}$ du catalogue avec 96\,Hz et sélectionner le premier modèle qui satisfait cette condition, sans surdimensionner inutilement le choix (coût, encombrement).
\end{tcolorbox}

\end{tcolorbox}

% -----------------------------------------------------------------------
%  4. ATTENTES DU JURY
% -----------------------------------------------------------------------

\begin{tcolorbox}[
  jury,
  title={4. Attentes du jury - Erreurs récurrentes et points de vigilance},
]

\begin{enumerate}[label=\textbf{\arabic*.}, leftmargin=*, itemsep=5pt]

  \item \textbf{Identifier le capteur depuis la référence.} Lire le datasheet fourni ; ne pas deviner. La confusion PT100 (passif, résistance) / thermocouple (actif, FEM de l'ordre du mV) est fréquente.

  \item \textbf{Justifier le choix chiffré.} Jury 2020 : "Le choix du capteur est rarement justifié." Toujours comparer avec les critères numériques (plage, tolérance, fréquence).

  \item \textbf{Pont de Wheatstone.} Jury 2025 : moins de 20\,\% des candidats traitent correctement. Mémoriser $V_{out} \approx (V_{cc}/4) \cdot (\Delta R/R_0)$ et vérifier $\Delta R \ll R_0$.

  \item \textbf{Signal 4-20\,mA.} $4\,\mathrm{mA}$ correspond à la valeur minimale, pas à zéro. La mise à l'échelle part de $I_s - 4$ et non de $I_s$.

  \item \textbf{Gain du conditionneur.} Calculer analytiquement : $G = (V_{max,CAN} - V_{min,CAN}) / (V_{max,capteur} - V_{min,capteur})$. Ne pas oublier l'offset si la plage capteur est décalée.

  \item \textbf{Relier capteur et CAN.} La tension de sortie du conditionneur doit \textbf{couvrir exactement} $[0\,;\,V_{ref}]$ - utiliser tout le codage du CAN pour maximiser la résolution effective.

\end{enumerate}

\medskip
\begin{tcolorbox}[
  colback=colJuryFrame!12!colJury,
  colframe=colJuryFrame,
  left=4pt, right=4pt, top=2pt, bottom=2pt,
  fontupper=\small\itshape
]
Rapport jury 2025 : «~L'intérêt d'un pont de Wheatstone et l'expression simplifiée de la tension de sortie ont été traités par moins de vingt pour cent des candidats.~»
\end{tcolorbox}

\end{tcolorbox}

% -----------------------------------------------------------------------
%  PIED - PRESENCE DANS LES SUJETS
% -----------------------------------------------------------------------

\begin{tcolorbox}[
  ficheSection=colPresence,
  title={Présence dans les sujets analysés (2013-2026)},
]
\begin{center}
\footnotesize\renewcommand{\arraystretch}{1.15}
\begin{tabular}{lccccccccccccc}
\toprule
\textbf{Année} & 13 & 14 & 15 & 16 & 17 & 18 & 19 & 20 & 21 & 22 & 23 & 24 & 25 \\
\midrule
E1 (commune)   & - & $\checkmark$ & - & - & $\checkmark$ & - & - & $\checkmark$ & $\checkmark$ & $\checkmark$ & - & - & $\checkmark$ \\
E2 discipl. IE & - & - & - & - & - & - & $\checkmark$ & - & - & $\checkmark$ & $\checkmark$ & $\checkmark$ & $\checkmark$ \\
E2 appl. comm. & - & - & - & - & - & - & - & - & - & - & - & - & $\checkmark$ \\
\bottomrule
\end{tabular}
\end{center}
\smallskip
\footnotesize\textbf{Sous-thèmes fréquents :}
PT100 ($R_T$, $\alpha$) - thermocouple (type J/K, FEM) - pont de Wheatstone - signal 4-20\,mA - montage AOP soustracteur - capteur inductif (fréquence) - mise à l'échelle CAN

\end{tcolorbox}

\CAPETENDFILE
